\section{Methods}

Every orchestration topology carries implicit assumptions about task structure. This section formalizes both sides of the alignment framework: the task-structure dimensions that characterize what a task demands, and the topology priors that characterize what each orchestration pattern assumes.

\subsection{Task-Structure Dimensions}\label{task-structure-dimensions}

We characterize tasks along three dimensions that capture the structural demands an orchestration topology must satisfy. Each is defined operationally so that annotators can assign scores from ground-truth trajectories.

\textbf{Decomposability (D).} Given a task $\tau$ with a reference solution of $n$ steps, let $k$ denote the number of steps whose inputs depend only on the original task specification. Then $D = k / n \in [0, 1]$. A coding task requiring bug localization across three independent modules followed by a unified patch exhibits $D \approx 0.6$: the localization steps are independent but patch synthesis depends on all of them.

\textbf{Iterativeness (I).} Let $r$ denote the number of steps that revise, override, or extend the output of a previously completed step. Then $I = r / n \in [0, 1]$. Reasoning tasks involving multi-hop logic score $I \approx 0.7$ because intermediate conclusions must be revised as new evidence accumulates across reasoning steps.

\textbf{Tool Diversity (T).} $T \in \{1, 2, \ldots, n\}$ counts the number of distinct tool categories the task requires, drawn from six canonical categories: code execution, web navigation, file I/O, search/retrieval, mathematical computation, and API invocation. A research task requiring web search, document reading, data extraction, and cross-source verification scores $T = 4$. We normalize $T$ to $[0, 1]$ by dividing by the maximum observed value across our task corpus, producing a unit cube for geometric computation.

\subsection{Topology Taxonomy}\label{topology-taxonomy}

Each orchestration topology encodes a structural prior: an implicit assumption about which region of D-I-T space the tasks it will face occupy. We characterize five topologies below. The first three are evaluated experimentally; the remaining two are positioned in D-I-T space for completeness but not yet evaluated.

\textbf{Flat conversation} organizes agents as equal participants in a multi-turn dialogue loop with shared tool access \cite{wu_2023}. No agent has authority over another; coordination emerges from conversational context. Prior: $(D_{\text{prior}}, I_{\text{prior}}, T_{\text{prior}}) = (0.2, 0.8, 0.3)$. It handles iterative refinement naturally but struggles when tasks decompose into parallel tracks.

When tasks decompose into parallel tracks, flat conversation breaks down --- but hierarchical SOP thrives. \textbf{Hierarchical SOP} assigns a planner agent to decompose tasks into subtask specifications, each executed by a specialist, with verification gates between stages \cite{hong_2023}. Prior: $(0.8, 0.2, 0.5)$. Its failure mode is the mirror image of flat's strength: tasks requiring iterative revision of earlier outputs resist clean decomposition, and the planner's subtask boundaries sever the dependencies that matter most.

\textbf{Debate} structures agents as adversarial reviewers who propose, critique, and refine solutions through multiple rounds of structured argumentation. Agents specialize in generating and stress-testing claims, surfacing errors that a single perspective would miss. Prior: $(0.3, 0.7, 0.3)$.

The framework extends to topologies we have not yet evaluated. \textbf{Role-playing} assigns distinct personas who collaborate through structured dialogue in fixed-turn exchanges \cite{li_2023}. Prior: $(0.5, 0.5, 0.3)$. \textbf{RL-orchestrated} topologies use a learned policy to dynamically assign agents based on runtime signals \cite{evolving_2025}. Because the policy adapts, RL topologies make weaker a priori assumptions, occupying a broad central region at $(0.5, 0.5, 0.5)$ with larger uncertainty bounds. Both are positioned in D-I-T space to generate testable predictions for future work.

\subsection{The Alignment Hypothesis Formalized}\label{the-alignment-hypothesis-formalized}

The topology-task alignment hypothesis states that for a given task $\tau$ with coordinates $(D_\tau, I_\tau, T_\tau)$, the topology whose structural prior $(D_p, I_p, T_p)$ minimizes the alignment distance

$$\delta(\tau, p) = \sqrt{w_D(D_\tau - D_p)^2 + w_I(I_\tau - I_p)^2 + w_T(T_\tau - T_p)^2}$$

will achieve the highest performance on $\tau$ among the evaluated topologies. Here $w_D, w_I, w_T \geq 0$ are dimension weights satisfying $w_D + w_I + w_T = 1$, estimated from a held-out validation set. The hypothesis is falsifiable: if topology effectiveness is unrelated to task structure, alignment distance will show no correlation with performance ranking.

Two predictions follow. First, no single topology should dominate across all task categories, because the three categories occupy different D-I-T regions. Second, tasks with extreme D-I-T profiles should show the largest performance gaps between best-matched and worst-matched topologies.

\subsection{OrchestraBench: Topology Implementations}\label{orchestrabench-topology-implementations}

OrchestraBench implements the three evaluated topologies (flat, hierarchical, debate) as controlled wrappers around a shared infrastructure stack. Each wrapper receives the same inputs (task specification, LLM backbone access, tool suite, token budget) and differs only in how it organizes agent interactions.

The shared infrastructure has four components. The \textbf{LLM backbone} provides all generation through a single endpoint (Claude Opus 4.6, with a secondary check on Claude Sonnet 4.6) with identical temperature and maximum token settings. The \textbf{tool suite} exposes all canonical tool categories through a unified interface (code interpreter, terminal, web browser, file system). The \textbf{context budget} uses the full 1M token context window, enforced at the wrapper level; topologies that spend tokens on inter-agent coordination reduce the budget available for task-relevant generation. The \textbf{evaluation harness} records all interactions, tool calls, intermediate outputs, and final answers in a standardized trace format.

Each wrapper handles agent instantiation, message routing, and termination. The wrappers are 200--600 lines of Python, kept deliberately lightweight so implementation quality does not confound the topology comparison.

\subsection{Task Sampling and Annotation}\label{task-sampling-and-annotation}

We evaluate on 82 custom tasks (12 easy, 70 hard) across three categories chosen for structural diversity. Tasks are hand-constructed to cover the full D-I-T space rather than sampled from existing benchmarks, giving us control over structural diversity at the cost of external validity (see Section~\ref{sec:results-limitations}).

\textbf{Coding} (26 tasks): implementation, debugging, and refactoring tasks. These cluster at high $D$ (mean 0.63), moderate $T$ (mean 2.3), and low $I$ (mean 0.27), because bug localization, patching, and testing are separable stages.

\textbf{Research} (30 tasks, including multi-domain tasks requiring cross-category tool use): information synthesis, cross-source verification, and tasks requiring both research and coding or reasoning components. These cluster at moderate-to-high $T$ (mean 3.7), moderate $D$ (mean 0.48), and moderate $I$ (mean 0.39) due to diverse tool requirements.

\textbf{Reasoning} (26 tasks): logic, game theory, and constraint satisfaction. Reasoning tasks proved hardest to annotate --- multi-hop logic required judgment calls about where revision begins and sequential extension ends. They cluster at high $I$ (mean 0.59), low $D$ (mean 0.24), and low $T$ (mean 1.4).

The authors independently score each task's $D$, $I$, and $T$ from reference solution trajectories. Inter-annotator agreement (Cohen's $\kappa$): 0.83 for D, 0.79 for I, 0.91 for T. Disagreements exceeding 0.2 on any dimension are resolved through discussion to consensus.

\subsection{Evaluation Protocol}\label{evaluation-protocol}

Each of the 82 tasks runs under all three evaluated topologies on the primary backbone (Claude Opus 4.6) with three different random seeds, producing 738 primary runs (82 $\times$ 3 $\times$ 3). A secondary backbone check runs all 12 easy tasks across all three topologies on Claude Sonnet 4.6 (36 runs). An agent-count control experiment gives a single flat agent 2$\times$ the token budget on 20 hard tasks (20 runs). Grand total: 794 runs. We measure three metrics: \textbf{task success rate} (binary correct/incorrect), \textbf{token efficiency} (total tokens consumed), and \textbf{wall-clock latency}. We report mean accuracy with 95\% confidence intervals from the three replicates. Alignment prediction accuracy is evaluated via a DIT routing classifier that selects the best topology based on alignment distance (Section~\ref{the-alignment-hypothesis-formalized}).
